%! TeX root = Realtorio.tex
\documentclass{article}
\usepackage[utf8]{inputenc}
\usepackage{amsmath}

\begin{document}

\section{Resolução da Questão 2}
\smallskip

No âmbito da resolução da questão, é pedido que seja calculado um valor aproximado da soma da série dada por:
\begin{equation}
    S = 18 \sum_{k=1}^{\infty} \frac{k!^2}{k^2 (2k)!}
\end{equation}
\bigskip

Inicialmente, procuramos definir uma simplificação para o termo geral da série S, dado por: $a_k = \frac{k!^2}{k^2 (2k)!}$, de modo a obter uma nova expressão através da qual possamos realizar o cálculo aproximado. Para isso, analisamos o termo seguinte, ou seja, $a_{k+1} = \frac{(k+1)!^2}{(k+1)^2 (2(k+1))!}$, e realizamos a divisão entre os dois termos, obtendo:

\begin{align*}
    \frac{a_{k+1}}{a_k} &= \frac{(k+1)!^2}{(k+1)^2 (2(k+1))!} \cdot \frac{k^2 (2k)!}{k!^2} \\
    &= \frac{(k+1)^2 k!^2}{(k+1)^2 (2k+2)(2k+1)(2k)!} \cdot \frac{k^2 (2k)!}{k!^2} \\
    &= \frac{k^2}{(2k+1)(2k+2)} \\
    &= \frac{k^2}{4k^2 + 6k + 2}
\end{align*}
\bigskip

Após esta simplificação, deduzimos que o termo geral da série S pode ser reescrito como:
\begin{equation}
    a_{k+1} = a_k \cdot \frac{k^2}{4k^2 + 6k + 2}
\end{equation}
\medskip

De modo a majorar o termo geral da série, podemos observar que, como k é um número natural e maior ou igual a 1, temos $4k^2 + 6k + 2 > 4k^2$, o que nos permite escrever a seguinte desigualdade:
\begin{equation}
    \frac{k^2}{4k^2 + 6k + 2} < \frac{k^2}{4k^2} = \frac{1}{4}
\end{equation}
\bigskip

Dessa forma, podemos concluir que, como S = $18 \sum_{k=1}^{\infty} a_k$ e k maior ou igual a 1, temos que:
\smallskip

\begin{equation}
    R_n < 18 \cdot a_k \sum_{j=0}^{\infty} \frac{1}{4^j} = 18 \cdot a_k \cdot \frac{1}{1 - 1/4} = 18 \cdot \frac{4}{3} \cdot a_k = 24 \cdot a_k
\end{equation}
\bigskip   


\end{document}
